% 本文件由 工程文件/scripts/make_refs.py 从原文 p.381--386 机械搬运，
% 仅修重音与转义；人名/刊名保持原文。
% 已人工校订：合并行末断字、清理多余空格、书名/刊名包 \emph、
% Unicode 数学字符转 LaTeX；原 p.386 中被并入 [T] 的 [Ve] 条已拆回独立条目。
%
% 第二轮校订（2026-11-21，对照原书 p.381--386 页面图像逐条核对）：
%   * 89 条的键名、次序、人名、标题、刊名/丛书名、卷期、年份、页码与底本逐字符一致；
%   * 底本排印缺陷（缩写后漏空格、vol 后漏点：Akman.A / I.Mirković / J.L. Verdier /
%     Trans.Amer.Math.Soc. / Amer.Math.Society / vol 176）已按通行写法补齐分隔；
%   * 底本自身的误拼（decompositiopns、algbras、O'Raifertaigh、Matematicae、géometrie、
%     local systems Preprint 等）一律照录，逐条加 [Tutor Note] 说明，不作订正；
%   * 底本对刊名/书名的斜体极不统一（同一刊名在不同条目中或斜或正，如 Duke Math. Jour.
%     在 [Hit87] 斜体、在 [La87][La88] 正体），译文沿用「刊名/书名一律 \emph」的既定体例，
%     不随底本逐条反复；此点已在 工程文件/校对/refs_校对.md 中说明。
% 另：本工程导言区目前并未定义 \bdsection（只有 \bdsec），为免本文件被 \input
% 时报 Undefined control sequence，这里做一次带保护的兜底定义；若导言区已定义则不生效。
\makeatletter
\@ifundefined{bdsection}{\NewDocumentCommand{\bdsection}{s o m}{\section*{#3}}}{}
\makeatother
\bdsection*{参考文献}
% 页眉：\bdsection* 走 \section*，不改 marks，若不显式设置，这 5 页的页眉会残留
% 上一节「10.1. Delta 函数」（章内 \bdsec 自带 \markboth，而 \bdsection 没有）。
\markboth{参考文献}{参考文献}

% 注：原行首选项 style=nextline 在 ctex 下会报
%   ! Argument of \enit@postlabel@i has an extra }.
% （enumitem 的 nextline 机制假定 \@item 未被改动，与 ctex 不兼容；已实测），
% 且原书 p.381--386 的 [键] 本与正文同行、只是悬挂缩进，故去掉该项，其余选项不动。
% 书目是纯西文，LaTeX 默认在句末句点后加额外空档（"vectorielle.  In:" 会拉宽），
% 底本排的是法式等距（\frenchspacing），故本环境内关掉句末附加空档。
\frenchspacing
\begin{description}[leftmargin=3.2em,itemsep=2pt]
\item[{[Ak]}] F. Akman. A characterization of the differential in semi-infinite cohomology. \emph{Journal of Algebra}, vol. 162 (1993), 194--209.
\item[{[Ar]}] M. Artin. Algebraization of formal moduli: II. Existence of modifications. \emph{Annals of Math.}, vol. 91 (1970), No. 1, 88--135.
\item[{[At]}] M. F. Atiyah. Riemann surfaces and spin structures. \emph{Ann. Sci. École Norm. Sup.}, ser. 4, vol. 4 (1971), 47--62.
\item[{[Ba63]}] H. Bass. Big projective modules are free. \emph{Illinois J. of Math.}, vol. 7 (1963), 24--31.
\item[{[Ba68]}] H. Bass. \emph{Algebraic K-theory}. Benjamin, New York, 1968.
\item[{[BB81]}] A. Beilinson and J. Bernstein. Localisation de $\mathfrak{g}$-modules. \emph{Comptes Rendus Acad. Sci. Paris}, vol. 292, série I (1981), 15--18.
\item[{[BB93]}] A. Beilinson and J. Bernstein. A proof of Jantzen conjectures. In: I. M. Gelfand Seminar, \emph{Advances in Soviet Math.}, vol. 16, part 1, 1--50. Amer. Math. Soc., Providence, 1993.
\item[{[BD94]}] A. Beilinson and V. Drinfeld. Affine Kac-Moody algebras and polydifferentials. \emph{International Mathematics Research Notices}, 1994, No. 1, 1--11.
\item[{[BD96]}] A. Beilinson and V. Drinfeld. Quantization of Hitchin's fibration and Langlands program. In: \emph{Algebraic and Geometric Methods in Mathematical Physics} (A. Boutet de Monvel and V. A. Marchenko, editors), \emph{Mathematical Physics Studies}, vol. 19, 3--7. Kluwer Academic Publishers, Netherlands, 1996.
% [Tutor Note: 原书 p.381 此条只写 "To appear"（当时尚未出版）；该书后以 Chiral Algebras（AMS Colloq. Publ. 51）行世。底本如此，保留原样。]
\item[{[BD]}] A. Beilinson and V. Drinfeld. \emph{Chiral algebras}. To appear.
\item[{[BS]}] A. Beilinson and V. Schechtman. Determinant bundles and Virasoro algebras. \emph{Commun. Math. Phys.}, vol. 118 (1988), 651--701.
\item[{[BL]}] J. Bernstein and V. Lunts. Localization for derived categories of $(\mathfrak{g},K)$-modules. \emph{Journal of Amer. Math. Soc.}, vol. 8 (1995), No. 4, 819--856.
\item[{[BLa94]}] A. Beauville and Y. Laszlo. Conformal blocks and generalized theta functions. \emph{Commun. Math. Phys.}, vol. 164 (1994), No. 2, 385--419.
\item[{[BLa95]}] A. Beauville and Y. Laszlo. Un lemme de descente. \emph{Comptes Rendus Acad. Sci. Paris}, vol. 320, série I (1995), 335--340.
\item[{[BLaSo]}] A. Beauville, Y. Laszlo, and C. Sorger. The Picard group of the moduli of $G$-bundles on a curve. Preprint alg-geom/9608002.
\item[{[BR]}] I.N. Bernstein and B.I. Rozenfel'd. Homogeneous spaces of infinite-dimensional Lie algebras and characteristic classes of foliations. \emph{Russian Math. Surveys}, vol. 28 (1973), no. 4, 107--142 (Russian original: \emph{Uspekhi Mat. Nauk}, vol. 28 (1973), no. 4, 103--138)
\item[{[Bourb59]}] N. Bourbaki. \emph{Algèbre}, chap. IX. Hermann, Paris, 1959.
\item[{[Bour68]}] N. Bourbaki. \emph{Groupes et algèbres de Lie}, chap. IV--VI. Hermann, Paris, 1968.
\item[{[Bour75]}] N. Bourbaki. \emph{Groupes et algèbres de Lie}, chap. VII, VIII. Hermann, Paris, 1975.
\item[{[De-Mi]}] P. Deligne and J. Milne. Tannakian categories. In: \emph{Hodge Cycles, Motives, and Shimura Varieties}, \emph{Lecture Notes in Mathematics}, vol. 900, 101--228. Springer-Verlag, 1982.
\item[{[Del70]}] P. Deligne. \emph{Équations différentielles à points singuliers réguliers}. \emph{Lecture Notes in Math}, vol. 163. Springer-Verlag, 1970.
\item[{[Del73]}] P. Deligne. La formule de dualité globale. In: \emph{SGA 4}, tome 3, \emph{Lecture Notes in Mathematics}, vol. 305, 481--587. Springer-Verlag, 1973.
\item[{[Del87]}] P. Deligne. Le déterminant de la cohomologie. In: \emph{Current Trends in arithmetical algebraic geometry}, \emph{Contemporary Mathematics}, vol. 67, 93--177. Amer. Math. Soc., Providence, 1987.
\item[{[Del91]}] P. Deligne. Catégories tannakiennes. In: \emph{Grothendieck Festschrift}, vol. II (P. Cartier et al., editors), 111--195, Birkhauser, Boston, 1991.
% [Tutor Note: 原书 p.382 书名印作 "La géometrie des groupes classiques"，géometrie 缺第二个重音（正字 géométrie）；底本如此，保留原样。]
\item[{[D71]}] J. Dieudonné. \emph{La géometrie des groupes classiques}. \emph{Ergebnisse der Mathematik und ihrer Grenzgebiete}, Band 5. Springer-Verlag, 1971.
\item[{[DS85]}] V.G. Drinfeld and V.V. Sokolov. Lie algebras and equations of Korteweg-deVries type. \emph{Journal of Soviet Math}, vol. 30 (1985), 1975--2036. Reprinted in [W], pp. 25--86.
\item[{[DSim]}] V. Drinfeld and C. Simpson. $B$-structures on $G$-bundles and local triviality. \emph{Mathematical Research Letters}, vol. 2 (1995), No. 6, 823--829.
\item[{[Fal93]}] G. Faltings. Stable $G$-bundles and projective connections. \emph{Journal of Algebraic Geometry}, vol. 2 (1993), No. 3, 507--568.
\item[{[Fal94]}] G. Faltings. A proof for the Verlinde formula. \emph{Journal of Algebraic Geometry}, vol. 3 (1994), No. 2, 347--374.
\item[{[F84]}] B.L. Feigin. The semi-infinite homology of Kac--Moody and Virasoro Lie algebras. \emph{Russian Math. Surveys}, vol. 39 (1984), no. 2, 155--156 (Russian original: \emph{Uspekhi Mat. Nauk}, vol. 39 (1984), no. 2, 195--196)
\item[{[F91]}] B. Feigin. Differential operators on the moduli space of $G$-bundles on algebraic curve and Lie algebra cohomologies. In: \emph{Special Functions}, editors M. Kashiwara, T. Miwa, 90--103. Springer-Verlag, Tokyo, 1991.
\item[{[FF92]}] B. Feigin and E. Frenkel. Affine Kac-Moody algebras at the critical level and Gelfand-Dikii algebras. In: \emph{Infinite Analysis}, editors A. Tsuchiya, T. Eguchi, M. Jimbo, \emph{Adv. Ser. in Math. Phys.}, vol. 16, 197--215. World Scientific, Singapore, 1992.
\item[{[Fr]}] E. Frenkel. Affine algebras, Langlands duality, and Bethe ansatz. In: \emph{Proceedings of the International Congress on Mathematical Physics}, Paris, 1994 (D. Iagolnitzer editor), 606--642. International Press, 1995.
\item[{[Fr81]}] I. B. Frenkel. Two constructions of affine Lie algebra representations and boson-fermion correspondence in quantum field theory. \emph{Journal of functional analysis}, vol. 44 (1981), 259--327.
\item[{[FGZ86]}] I. B. Frenkel, H. Garland, and G. J. Zuckerman. Semi-infinite cohomology and string theory. \emph{Proc. Natl. Acad. Sci. USA} 83 (1986), pp. 8442--8446.
\item[{[Ga2001]}] D. Gaitsgory. Construction of central elements in the affine Hecke algebra via nearby cycles. \emph{Invent. Math.}, vol. 144 (2001), no. 2, 253--280.
\item[{[GD76]}] I.M. Gelfand and L.A. Dikii. Fractional powers of operators and Hamiltonian systems. \emph{Functional Analysis and Applications}, vol. 10 (1976), no. 4, 259--273 (Russian original: pp. 13--29).
\item[{[GD78]}] I.M. Gelfand and L.A. Dikii. A family of Hamiltonian structures connected with integrable nonlinear differential equations. Preprint Inst. Appl. Math. no. 136 (Moscow, 1978). Reprinted in [W], pp. 3-24.
\item[{[Gr64]}] A. Grothendieck. \emph{Éléments de géométrie algébrique}. IV. Étude locale des schémas et de morphismes de schémas (première partie). \emph{Publications Mathématiques IHES}, No. 20, 1964.
\item[{[Gr67]}] A. Grothendieck. \emph{Éléments de géométrie algébrique}. IV. Étude locale des schémas et de morphismes de schémas (quatrième partie). \emph{Publications Mathématiques IHES}, No. 32, 1967.
\item[{[Gr68]}] A. Grothendieck. Crystals and the De Rham cohomology of schemes. In: \emph{Dix Exposés sur la Cohomologie des Schémas}, 306--358. North-Holland, Amsterdam, 1968.
\item[{[Gr68a]}] A. Grothendieck. \emph{Catégories cofibrées additives et complexe cotangent relatif}. \emph{Lecture Notes in Math.}, vol. 79. Springer-Verlag, 1968.
\item[{[Gi87]}] V. Ginzburg. Equivariant cohomology and Kähler geometry. \emph{Functional Anal. Appl.}, vol. 21 (1987), no. 4, 271--283 (Russian original: pp. 19--34).
\item[{[Gi95]}] V. Ginzburg. Perverse sheaves on a loop group and Langlands' duality. Preprint alg-geom/9511007.
\item[{[Gi97]}] V. Ginzburg. Global nilpotent variety is Lagrangian. Preprint alg-geom/9704005.
\item[{[GK]}] I.M. Gel'fand and D.A. Kazhdan. Some problems of differential geometry and the calculation of cohomologies of Lie algebras of vector fields. \emph{Soviet Math. Doklady}, vol. 12 (1971), no. 5, 1367--1370 (Russian original: \emph{Doklady Akademii Nauk SSSR}, vol. 200 (1971), no. 2, 269--272).
\item[{[GKF]}] I.M. Gel'fand, D.A. Kazhdan, and D.B. Fuks. The actions of infinite-dimensional Lie algebras. \emph{Functional Anal. Appl.}, vol. 6 (1972), no. 1, 9--13 (Russian original: pp. 10--15).
\item[{[Ha88]}] T. Hayashi. Sugawara operators and Kac-Kazhdan conjecture. \emph{Inventiones Math.}, vol. 94 (1988), 13--52.
\item[{[Hit87]}] N. Hitchin. Stable bundles and integrable systems. \emph{Duke Math. Jour.}, vol. 54 (1987), 91--114.
% [Tutor Note: 原书 p.384 标题印作 "Bruhat decompositiopns"，应为 decompositions；底本如此，保留原样。]
\item[{[IM]}] N. Iwahori and H. Matsumoto. On some Bruhat decompositiopns and the structure of the Hecke rings of $\mathfrak{p}$-adic Chevalley groups. \emph{Publications Mathématiques IHES}, No. 25, 1965.
\item[{[J]}] C.U. Jensen. On the vanishing of $\lim\limits_{\longleftarrow}{}^{(i)}$. \emph{Jour. of Algebra}, vol. 15 (1970), 151--166.
\item[{[Kac90]}] V.G. Kac. \emph{Infinite dimensional Lie algebras}. Cambridge University Press, 1990.
% [Tutor Note: 原书 p.384 出版项印作 "Amer.Math.Society"，缩写点后无空格；底本如此，译文按通行写法排作 "Amer. Math. Society"。]
\item[{[Kac97]}] V.G. Kac. \emph{Vertex algebras for beginners}. \emph{University lecture series}, vol. 10. Amer. Math. Society, Providence, USA, 1997.
\item[{[KK]}] V.G. Kac and D.A. Kazhdan. Structure of representations with highest weight of infinite-dimensional Lie algebras. \emph{Adv. in Math.}, vol. 34 (1979), No. 1, 97--108.
\item[{[KP]}] V.G. Kac and D.H. Peterson. Spin and wedge representations of infinite-dimensional Lie algebras and groups. \emph{Proc. Nat. Acad. Sci. USA}, vol. 78 (1981), No. 6, 3308--3312.
\item[{[KR]}] V.G. Kac and A.K. Raina. \emph{Bombay lectures on highest weight representations of infinite dimensional Lie algebras}. World Scientific, Singapore, 1987.
\item[{[Kap91]}] M. Kapranov. On DG-modules over the de Rham complex and the vanishing cycles functor. In: \emph{Lecture Notes in Mathematics}, vol. 1479, 57--86. Springer-Verlag, Berlin, 1991.
% [Tutor Note: 原书 p.384 此处 "local systems" 后面漏排句点（底本如此）；译文保留缺失，未代补。]
\item[{[Kap97]}] M. Kapranov. Injective resolutions of $BG$ and derived moduli spaces of local systems Preprint alg-geom/9710027.
\item[{[Kn-Mu]}] F. Knudsen and D. Mumford. The projectivity of the moduli space of stable curves I: preliminaries on ``det'' and ``div''. \emph{Math. Scand.}, vol. 39 (1976), 19--55.
\item[{[Kas]}] M. Kashiwara. Representation theory and $D$-modules on flag varieties. \emph{Astérisque}, vol. 173--174 (1989), 55--109.
\item[{[KT95]}] M. Kashiwara and T. Tanisaki. Kazhdan-Lusztig conjecture for affine Lie algebras with negative level. \emph{Duke Math. Jour.}, vol. 77 (1995), No. 1, 21--62.
\item[{[Ko61]}] B. Kostant. Lie algebra cohomology and the generalized Borel-Weil theorem. \emph{Annals of Math.}, vol. 74 (1961), No. 2, 329--387.
\item[{[Ko63]}] B. Kostant. Lie group representations on polynomial rings. \emph{Amer. Jour. of Math.}, vol. 85 (1963), 327--404.
\item[{[Ko78]}] B. Kostant. On Whittaker vectors and representation theory. \emph{Inventiones Math.}, vol. 48 (1978), 101--184.
\item[{[KS]}] B. Kostant and S. Sternberg. Symplectic reduction, BRS cohomology, and infinite-dimensional Clifford algebras. \emph{Annals of Physics}, vol. 176 (1987), 49--113.
\item[{[La-So]}] Y. Laszlo and C. Sorger. The line bundles on the moduli of parabolic $G$-bundles over curves and their sections. \emph{Ann. Sci. École Norm. Sup.}, ser. 4, vol. 30 (1997), 499--525.
\item[{[La87]}] G. Laumon. Correspondance de Langlands géométrique pour les corps de fonctions. \emph{Duke Math. Jour.}, vol. 54 (1987), 309--359.
\item[{[La88]}] G. Laumon. Un analogue global du cône nilpotent. \emph{Duke Math. Jour.}, vol. 57 (1988), 647--671.
\item[{[LMB93]}] G. Laumon and L. Moret-Bailly. \emph{Champs Algébriques}. Preprint 92-42, Université Paris-Sud, 1992.
\item[{[Lu83]}] G. Lusztig. Singularities, character formulas, and a $q$-analog of weight multiplicities. \emph{Astérisque}, vol. 101-102 (1983), 208-229.
\item[{[Ma87]}] K. Mackenzie. \emph{Lie Groupoids and Lie Algebroids in Differential Geometry}. \emph{London Math. Soc. Lecture Note Ser.}, vol. 124. Cambridge Univ. Press, 1987.
\item[{[Ma96]}] K. Mackenzie. Lie algebroids and Lie pseudoalgebras. \emph{Bull. London Math. Soc.}, vol. 27 (1995), part 2, 97--147.
\item[{[Mu]}] D. Mumford. Theta characteristics on an algebraic curve. \emph{Ann. Sci. École Norm. Sup.}, ser. 4, vol. 4 (1971), 181--192.
\item[{[MV]}] I. Mirković and K. Vilonen. Perverse sheaves on Loop Grassmannians and Langlands duality. Preprint alg-geom/9703010.
\item[{[Pa]}] B. Pareigis. Non-additive ring and module theory. II. $C$-categories, $C$-functors and $C$-morphisms. \emph{Publ. Math. Debrecen}, vol. 24 (1977), No. 3--4, 351--361.
\item[{[Pa76]}] A. Parshin. On the arithmetic of two-dimensional schemes, I. Distributions and residues. \emph{Mathematics of USSR -- Izvestiya}, vol. 10 (1976), 695--729 (Russian original: vol. 40, 736--773).
% [Tutor Note: 原书 p.386 作者名印作 "O'Raifertaigh"，正字为 O'Raifeartaigh；底本如此，保留原样。]
\item[{[Phys]}] L. Fehér, L. O'Raifertaigh, and I. Tsutsui. The vacuum preserving Lie algebra of a classical $W$-algebra. \emph{Phys. Lett. B} 316 (1993), 275--281.
\item[{[RG]}] M. Raynaud and L. Gruson. Critères de platitude et de projectivité. \emph{Inventiones Math.}, vol. 13 (1971), 1--89.
% [Tutor Note: 原书 p.386 标题印作 "Lie algbras"（应为 algebras），刊名印作 "Trans.Amer.Math.Soc."（缩写点后无空格）；误拼照录，空格按通行写法补齐。]
\item[{[R]}] G. S. Rinehart. Differential forms on general Lie algbras. \emph{Trans. Amer. Math. Soc.}, vol. 108 (1963), 195--222.
\item[{[Sa89]}] M. Saito. Induced $D$-modules and differential complexes. \emph{Bull. Soc. Math. France}, tome 117 (1989), fasc. 3, 361--387.
\item[{[Sa91]}] M. Saito. $D$-modules on analytic spaces. \emph{Publ. Res. Inst. Math. Sci.}, vol. 27 (1991), No. 2, 291--332.
\item[{[Se]}] J.-P. Serre. Modules projectifs et espaces fibrés à fibre vectorielle. In: \emph{Collected papers by J.-P. Serre}, vol. I, 531-543, Springer-Verlag, Berlin, 1986.
\item[{[T]}] J. Tate. Residues of differentials on curves. \emph{Ann. Sci. École Norm. Sup.}, ser. 4, vol. 1 (1968), 149--159.
\item[{[Ve]}] J. L. Verdier. Catégories dérivées. In: \emph{SGA 4$\frac{1}{2}$}, \emph{Lecture Notes in Mathematics}, vol. 569, 262--311. Springer-Verlag, Berlin, 1977.
\item[{[W]}] \emph{$W$-symmetry} (P. Bouwknegt and K. Schoutens editors). \emph{Adv. Ser. in Math. Phys.}, vol. 22. World Scientific, Singapore, 1995.
% [Tutor Note: 原书 p.386 刊名印作 "Expositiones Matematicae"，该刊实作 Expositiones Mathematicae；底本如此，保留原样。]
\item[{[We83]}] A. Weinstein. Sophus Lie and symplectic geometry. \emph{Expositiones Matematicae}, vol. 1 (1983), 95--96.
\item[{[We88]}] A. Weinstein. Coisotropic calculus and Poisson groupoids. \emph{Jour. Math. Soc. Japan}, vol. 40 (1988), No. 4, 705--727.
\item[{[We]}] A. Weinstein. Groupoids: unifying internal and external symmetry. \emph{Notices of AMS}, vol. 43 (1996), No. 7, 744--752.
\item[{[Yet]}] D. N. Yetter. Tangles in prisms, tangles in cobordisms. In: \emph{Topology '90} (B. Apanasov et al., editors), 399--443. Walter de Gruyter, Berlin, New York, 1992.
\end{description}
\nonfrenchspacing
